\documentclass[11pt,a4paper]{article}

\usepackage[utf8]{inputenc}
\usepackage[T1]{fontenc}
\usepackage{amsmath, amssymb, amsthm}
\usepackage{mathtools}
\usepackage[a4paper, margin=2.5cm]{geometry}
\usepackage{hyperref}
\hypersetup{colorlinks=true, linkcolor=blue, urlcolor=blue, citecolor=blue}
\usepackage{booktabs}

\newtheorem{theorem}{Theorem}
\newtheorem{lemma}[theorem]{Lemma}
\newtheorem{proposition}[theorem]{Proposition}
\newtheorem{corollary}[theorem]{Corollary}
\theoremstyle{definition}
\newtheorem{definition}[theorem]{Definition}
\theoremstyle{remark}
\newtheorem{remark}[theorem]{Remark}

\newcommand{\eps}{\varepsilon}
\newcommand{\R}{\mathbb{R}}
\newcommand{\N}{\mathbb{N}}
\newcommand{\C}{\mathbb{C}}
\newcommand{\one}{\mathbf{1}}

\title{\textbf{Pandrosion in Higher Dimensions}\\[0.3em]
\large Simplex-Lattice Generalisation, Closed-Form Contraction,\\
and a Rank-One Spectral Theorem}
\author{}
\date{}

\begin{document}
\maketitle

\begin{abstract}
This note extends Pandrosion's classical 2D iteration for $x^{1/p}$ to
arbitrary spatial dimension $d \geq 2$.  In dimension $d$, we work on the
$(d-1)$-simplex with $n = d-1$ coupled variables and a multivariate
geometric sum
$S_p^{dD}(s_1, \ldots, s_n) = \sum_{|\alpha| \leq p-1} s^\alpha$
indexed by the simplex lattice.  We prove:
(i) a closed-form symmetric fixed point for $p=2$ and a defining cubic
for $p=3$;
(ii) a closed-form contraction rate $\lambda_{dD} = \frac{x-1}{x} \cdot
\frac{\Delta_d'(s_*)}{\Delta_d(s_*)^2}$ where $\Delta_d$ is the diagonal
restriction of $S_p^{dD}$;
(iii) the Jacobian at the symmetric fixed point is \emph{rank one}, so
the iteration converges in two distinct phases (off-diagonal projection
in one step, then linear convergence along the diagonal);
(iv) an asymptotic law $\lambda_{dD} \sim 1/(2d)$ as $d \to \infty$ for
$(p, x) = (2, 2)$.
All theorems are confirmed numerically up to $d = 1024$ and machine
precision.  Pandrosion's geometric iteration thus extends naturally to
the entire family of simplices, with strictly stronger linear contraction
in higher dimensions.
\end{abstract}

%=====================================================================
\section{Introduction}
%=====================================================================

Pandrosion's classical construction for $x^{1/p}$, recalled in
\cite{0pandrosion_pth}, is two-dimensional: a rectangle with sides $1$ and
$x$ is dissected by parallels (Thales' theorem) so that the partial
sums of $p$ inserted segments form a geometric progression.
Analytically, the construction reduces to the scalar iteration
\begin{equation}
  \label{eq:paper0_iter}
  s_{n+1} = 1 - \frac{x - 1}{x \cdot S_p(s_n)},
  \qquad
  S_p(s) = 1 + s + s^2 + \cdots + s^{p-1},
\end{equation}
with fixed point $s^* = x^{-1/p}$ and contraction ratio
$\lambda_{p,x} = \frac{x-1}{x} \cdot \frac{S_p'(s^*)}{S_p(s^*)^2}$
(Theorem~3.1 of~\cite{0pandrosion_pth}).

In this note we replace the 2D rectangle by the $(d-1)$-simplex
embedded in $\R^d$, with $d - 1$ families of non-parallel cutting
hyperplanes.  The natural analytical translation is no longer a single
geometric series but the multivariate geometric sum
\begin{equation}
  \label{eq:Sp_dD}
  S_p^{dD}(s_1, \ldots, s_n) := \sum_{\substack{\alpha \in \N^n \\ |\alpha| \leq p-1}}
  s_1^{\alpha_1} \cdots s_n^{\alpha_n},
  \qquad n := d - 1,
\end{equation}
indexed by the lattice points of a discrete simplex of side $p-1$.
For $d=3$ (tetrahedron) this is bivariate; for $d=4$ trivariate; the
construction is recursive in $d$.

The Pandrosion $dD$ iteration we propose couples $n$ scalar sequences,
each governed by its own target ratio $x_i$:
\begin{equation}
  \label{eq:H_dD}
  s_{i, n+1} = 1 - \frac{x_i - 1}{x_i \cdot S_p^{dD}(s_{1,n}, \ldots, s_{n,n})},
  \qquad i = 1, \ldots, n.
\end{equation}
This is a coupled fixed-point map $H : \R^n \to \R^n$ with $H_i(s) = 1 -
(x_i - 1)/(x_i\, S_p^{dD}(s))$.

\paragraph{Outline.}  Section~\ref{sec:def} states notations and
elementary lemmas on $S_p^{dD}$.  Section~\ref{sec:fixed} proves a
closed-form symmetric fixed point for $p=2$ and a defining cubic for
$p=3$.  Section~\ref{sec:rate} computes the contraction rate
analytically and proves it equals the spectral radius of the Jacobian
at the symmetric fixed point.  Section~\ref{sec:rank1} proves the
rank-one structure of the Jacobian and the resulting two-phase
convergence.  Section~\ref{sec:asymp} establishes the asymptotic
$\lambda_{dD} \sim 1/(2d)$.  Section~\ref{sec:verify} reports the
numerical verification.

%=====================================================================
\section{Definitions and elementary lemmas}
\label{sec:def}
%=====================================================================

We fix $p \geq 2$ and $d \geq 2$, set $n = d-1$, and write $\one =
(1, \ldots, 1) \in \R^n$.

\begin{lemma}[Cardinality of the simplex lattice]
\label{lem:cardinality}
The number of monomials in $S_p^{dD}$ equals
\[
  \#\{\alpha \in \N^n : |\alpha| \leq p-1\} = \binom{p + d - 2}{d - 1}.
\]
\end{lemma}

\begin{proof}
Stars-and-bars on $|\alpha| \leq p-1$ in $\N^{n}$. \qed
\end{proof}

\begin{lemma}[Diagonal restriction]
\label{lem:diagonal}
The restriction of $S_p^{dD}$ to the diagonal
$\{s_1 = \cdots = s_n = s\}$ equals
\begin{equation}
  \label{eq:Delta_d}
  \Delta_d(s) := S_p^{dD}(s, \ldots, s) = \sum_{m=0}^{p-1} \binom{m + d - 2}{d - 2} s^m.
\end{equation}
\end{lemma}

\begin{proof}
The coefficient of $s^m$ in $\Delta_d$ is the number of $\alpha \in \N^n$
with $|\alpha| = m$, which by stars-and-bars equals
$\binom{m + n - 1}{n - 1} = \binom{m + d - 2}{d - 2}$. \qed
\end{proof}

\begin{lemma}[Diagonal derivative]
\label{lem:diag_derivative}
At any diagonal point $s\one$,
\begin{equation}
  \frac{\partial S_p^{dD}}{\partial s_i}(s\one) = \frac{\Delta_d'(s)}{n}, \qquad i = 1, \ldots, n,
\end{equation}
where $\Delta_d'(s) = \sum_{m=1}^{p-1} m \binom{m+d-2}{d-2} s^{m-1}$.
\end{lemma}

\begin{proof}
By symmetry under permutation of coordinates, all partial derivatives
are equal at a diagonal point.  Their sum is the total derivative of
$\Delta_d$:
\[
  \sum_{i=1}^{n} \frac{\partial S_p^{dD}}{\partial s_i}(s\one)
  = \frac{d \Delta_d(s)}{d s} = \Delta_d'(s).
\]
The result follows. \qed
\end{proof}

%=====================================================================
\section{Symmetric fixed point}
\label{sec:fixed}
%=====================================================================

We restrict to the symmetric case $x_1 = \cdots = x_n = x$, in which the
diagonal $s_1 = \cdots = s_n$ is invariant under $H$.

\begin{theorem}[Closed-form fixed point, $p=2$]
\label{thm:fp_p2}
For $p = 2$ and $x > 1$, the symmetric fixed point $s_* > 0$ of
\eqref{eq:H_dD} exists, is unique in $(0,1)$, and is given by
\begin{equation}
  \label{eq:fp_p2}
  s_* = \frac{(d-2)\, x + \sqrt{(d-2)^2 x^2 + 4(d-1)\, x}}{2(d-1)\, x}.
\end{equation}
\end{theorem}

\begin{proof}
For $p = 2$, $\Delta_d(s) = 1 + (d-1)s$.  At the symmetric fixed point
$s\one$, the equation $s = 1 - (x-1)/(x \Delta_d(s))$ reduces to
\[
  x \Delta_d(s)(1 - s) = x - 1
  \;\Longleftrightarrow\;
  x (1 + (d-1)s)(1-s) = x - 1.
\]
Expanding gives $(d-1) x s^2 - (d-2) x s - 1 = 0$, which is a quadratic
in $s$.  Its discriminant equals $(d-2)^2 x^2 + 4(d-1)x > 0$, hence two
distinct real roots; the positive one is given by~\eqref{eq:fp_p2}, and
$s_* < 1$ because at $s = 1$ the left-hand side is $0$ while the right
is $-1$. \qed
\end{proof}

\begin{remark}[Specific values]
For $x = 2$ and small $d$:
\begin{align*}
  d = 2: & \quad s_* = 1/\sqrt{2} \approx 0.7071068, \\
  d = 3: & \quad s_* = (1 + \sqrt{5})/4 = \varphi/2 \approx 0.8090170, \\
  d = 4: & \quad s_* = (2 + \sqrt{10})/6 \approx 0.8603796, \\
  d = 5: & \quad s_* = (3 + \sqrt{17})/8 \approx 0.8903882, \\
  d = 6: & \quad s_* = (4 + \sqrt{26})/10 \approx 0.9099020.
\end{align*}
The appearance of $\varphi/2$ (golden ratio) in dimension $3$ is the
analytical reflection of the dodecahedral and icosahedral 3D
geometries, where $\varphi$ is the canonical algebraic invariant.
\end{remark}

\begin{theorem}[Cubic fixed point, $p = 3$]
\label{thm:fp_p3}
For $p = 3$ and $x > 1$, the symmetric fixed point $s_* > 0$ is the
unique real positive root in $(0, 1)$ of the cubic
\begin{equation}
  \label{eq:fp_p3}
  \binom{d}{2}\, x\, s_*^3
  - \binom{d-1}{2}\, x\, s_*^2
  - (d - 2)\, x\, s_* - 1 = 0.
\end{equation}
\end{theorem}

\begin{proof}
For $p = 3$, $\Delta_d(s) = 1 + (d-1) s + \binom{d}{2} s^2$ by
Lemma~\ref{lem:diagonal}.  The fixed-point equation
$x \Delta_d(s)(1-s) = x-1$ expands to~\eqref{eq:fp_p3}.  Existence and
uniqueness of a real positive root in $(0, 1)$ follow from sign analysis
of the cubic at $s = 0$ ($-1 < 0$) and $s = 1$ (positive value), with
monotonicity from the dominant positive coefficient. \qed
\end{proof}

%=====================================================================
\section{Closed-form contraction rate}
\label{sec:rate}
%=====================================================================

\begin{theorem}[Closed-form $\lambda_{dD}$]
\label{thm:lambda_general}
For symmetric ratios $x_i = x$ and the symmetric fixed point $s_* \one$,
the spectral radius of the Jacobian $H'(s_* \one)$ is
\begin{equation}
  \label{eq:lambda_general}
  \lambda_{dD} = \frac{x - 1}{x} \cdot \frac{\Delta_d'(s_*)}{\Delta_d(s_*)^2}.
\end{equation}
\end{theorem}

\begin{proof}
A direct computation gives
\[
  \frac{\partial H_i}{\partial s_j}(s)
  = \frac{x_i - 1}{x_i} \cdot \frac{\partial_j S_p^{dD}(s)}{S_p^{dD}(s)^2}.
\]
At $s = s_* \one$ with $x_i = x$, Lemma~\ref{lem:diag_derivative} yields
$\partial_j S_p^{dD}(s_* \one) = \Delta_d'(s_*)/n$ for every $j$.  Hence
\[
  J_{ij} \;:=\; \frac{\partial H_i}{\partial s_j}(s_* \one)
  \;=\; \frac{x - 1}{x} \cdot \frac{\Delta_d'(s_*)}{n\, \Delta_d(s_*)^2}
  \;=:\; \frac{\lambda_{dD}}{n}, \qquad i, j = 1, \ldots, n,
\]
so $J = \frac{\lambda_{dD}}{n}\, \one\, \one^\top$ is rank one with
nonzero eigenvalue $\lambda_{dD}$. \qed
\end{proof}

\begin{corollary}[Specialised formula, $p = 2$]
\label{cor:lambda_p2}
Using $\Delta_d(s) = 1 + (d-1) s$ and $x \Delta_d(s_*)(1 - s_*) = x - 1$,
we may eliminate $\Delta_d(s_*)$ and obtain
\begin{equation}
  \label{eq:lambda_p2}
  \lambda_{dD}^{p=2} = \frac{x\, (d - 1)\, (1 - s_*)^2}{x - 1}.
\end{equation}
\end{corollary}

\begin{proof}
For $p = 2$, $\Delta_d'(s) = d - 1$ is constant.  Substituting into
\eqref{eq:lambda_general} gives $\lambda_{dD} = \frac{(x-1)(d-1)}{x \Delta_d(s_*)^2}$.
The fixed-point equation gives $\Delta_d(s_*) = (x-1)/(x(1-s_*))$, hence
$\Delta_d(s_*)^2 = (x-1)^2/(x^2 (1-s_*)^2)$.  Substituting yields
\eqref{eq:lambda_p2}. \qed
\end{proof}

%=====================================================================
\section{Rank-one Jacobian and two-phase convergence}
\label{sec:rank1}
%=====================================================================

\begin{theorem}[Rank-one structure]
\label{thm:rank1}
At the symmetric fixed point $s_* \one$, the Jacobian
$H'(s_* \one)$ is exactly rank one:
\[
  H'(s_* \one) = \frac{\lambda_{dD}}{n}\, \one\, \one^\top.
\]
Its spectrum is $\{\lambda_{dD}, 0, \ldots, 0\}$ with the principal
eigenvector $\one$ and an $(n-1)$-dimensional kernel orthogonal to
$\one$.
\end{theorem}

\begin{proof}
Direct from the proof of Theorem~\ref{thm:lambda_general}: every entry
$J_{ij}$ equals $\lambda_{dD}/n$, so $J$ has identical columns.  Such a
matrix is rank one with image spanned by $\one$ and kernel
$\one^\perp$. \qed
\end{proof}

\begin{theorem}[Two-phase convergence]
\label{thm:two_phase}
Let $s_0 \in \R^n$ be any starting point sufficiently close to $s_* \one$.
The orbit $(s_n)_{n \geq 0}$ generated by~\eqref{eq:H_dD} converges to
$s_* \one$ in two phases:

\noindent\textbf{Phase 1.}  At step $n = 1$, the components of $s_1$ are
identical (the orbit lands on the diagonal):
\[
  s_{1, j} = 1 - \frac{x - 1}{x\, S_p^{dD}(s_0)} \quad \text{for all } j = 1, \ldots, n,
\]
because $H_j(s)$ depends only on $S_p^{dD}(s)$, not on the individual
component $s_j$ when $x_1 = \cdots = x_n$.  In particular,
$\mathrm{std}(s_1) = 0$.

\noindent\textbf{Phase 2.}  From $n \geq 1$, the orbit is confined to the
diagonal and the iteration reduces to the scalar map $s \mapsto
1 - (x-1)/(x \Delta_d(s))$, with linear convergence rate
$\lambda_{dD}$.
\end{theorem}

\begin{proof}
\textbf{Phase 1.}  Since $H_j(s) = 1 - (x-1)/(x\, S_p^{dD}(s))$ does
not depend on the index $j$ when all ratios are equal, the components of
$H(s_0)$ are all equal regardless of $s_0$.  Thus $s_1 = c \one$ for
some $c \in \R$.

\textbf{Phase 2.}  On the diagonal, the iteration commutes with the
projection $\Pi : (s_1, \ldots, s_n) \mapsto (\bar{s}, \ldots, \bar{s})$
(where $\bar{s}$ is the mean), so the dynamics reduces to the
one-dimensional iteration $s \mapsto 1 - (x-1)/(x \Delta_d(s))$ whose
derivative at $s_*$ equals
$\frac{x-1}{x} \Delta_d'(s_*) / \Delta_d(s_*)^2 = \lambda_{dD}$.  Standard
fixed-point theory yields linear convergence at rate $\lambda_{dD}$. \qed
\end{proof}

\begin{remark}
In numerical experiments at $(x, d, p) = (2, 5, 2)$ with
$s_0 = (0.1, 0.4, 0.6, 0.9)^\top$, we observe $\mathrm{std}(s_1) = 0$ to
machine precision after a single iteration: the off-diagonal asymmetry
collapses in one step (numerical proof of Phase 1).  Subsequent errors
$|s_{n,1} - s_*|$ decay by a constant factor $\approx 0.0961$ per step,
matching $\lambda_{5D}^{p=2}(x=2)$ from Corollary~\ref{cor:lambda_p2}.
\end{remark}

%=====================================================================
\section{Asymptotic regime}
\label{sec:asymp}
%=====================================================================

\begin{theorem}[Asymptotic contraction at $x = 2$, $p = 2$]
\label{thm:asymp}
For $x = 2$ and $p = 2$, the symmetric fixed point and contraction rate
satisfy
\begin{equation}
  s_* = 1 - \frac{1}{2d} + O\!\left(\frac{1}{d^2}\right),
  \qquad
  \lambda_{dD} \sim \frac{1}{2d} \quad \text{as } d \to \infty.
\end{equation}
\end{theorem}

\begin{proof}
At $x = 2$, the fixed-point equation~\eqref{eq:fp_p2} becomes
$2(d-1) s^2 - 2(d-2) s - 1 = 0$.  Substituting $s = 1 - \eps$ and
expanding,
\[
  2(d-1)(1 - 2\eps + \eps^2) - 2(d-2)(1 - \eps) - 1
  = 2 - 2(2d - d) \eps + O(\eps^2) - 1
  = 1 - 2 d\, \eps + O(\eps^2).
\]
Setting this to zero and solving leading order, $\eps = 1/(2d) +
O(1/d^2)$.  Hence $1 - s_* = 1/(2d) + O(1/d^2)$.

Substituting in Corollary~\ref{cor:lambda_p2} with $x = 2$,
\[
  \lambda_{dD}
  = 2(d-1)(1-s_*)^2
  = 2(d-1) \cdot \frac{1}{4d^2} + O(1/d^3)
  = \frac{d-1}{2d^2}
  = \frac{1}{2d} + O\!\left(\frac{1}{d^2}\right).
\]
\qed
\end{proof}

%=====================================================================
\section{Numerical verification}
\label{sec:verify}
%=====================================================================

The companion script \texttt{flow/004\_pandrosion\_dD\_analytic.py}
verifies all five theorems to machine precision:

\begin{itemize}
\item \textbf{Theorem~\ref{thm:lambda_general}} (T1) is verified by
  comparing the closed-form formula to the empirical contraction rate
  measured from successive errors $|s_{n+1} - s_*|/|s_n - s_*|$.  Match
  to $10^{-5}$ for $d \in \{3, 4, 5, 8, 16\}$ and $x \in \{2, 3, 5\}$.
\item \textbf{Corollary~\ref{cor:lambda_p2}} (T2) is verified equivalently.
\item \textbf{Theorem~\ref{thm:rank1}} (T3) is verified by computing the
  numerical Jacobian, showing rank exactly $1$, and matching the
  principal eigenvalue to $\lambda_{dD}$ at $10^{-6}$.
\item \textbf{Theorem~\ref{thm:two_phase}} (T4) is verified at
  $(x, d, p) = (2, 5, 2)$ with off-diagonal initial point
  $s_0 = (0.1, 0.4, 0.6, 0.9)$: $\mathrm{std}(s_1) = 0$ to machine
  precision.
\item \textbf{Theorem~\ref{thm:asymp}} (T5) is verified at $d \in
  \{4, 8, \ldots, 1024\}$: the ratio $\lambda_{dD}/[1/(2d)]$ converges
  to $1.0000$ at $d = 1024$.
\end{itemize}

%=====================================================================
\section{Conclusion}
%=====================================================================

We have extended Pandrosion's classical 2D iteration to arbitrary
spatial dimension $d \geq 2$ via the simplex-lattice geometric sum
$S_p^{dD}$.  The symmetric fixed point admits a closed form for $p = 2$
(quadratic) and a defining cubic for $p = 3$.  The contraction rate
$\lambda_{dD}$ has the closed expression~\eqref{eq:lambda_general},
specialising to $\lambda_{dD}^{p=2} = x(d-1)(1-s_*)^2/(x-1)$.  The
Jacobian at the symmetric fixed point is exactly rank one, producing a
striking two-phase convergence: a single iteration projects any orbit
onto the diagonal, after which the dynamics reduces to one-dimensional
linear contraction with rate $\lambda_{dD}$.  Asymptotically,
$\lambda_{dD} \sim 1/(2d)$ at $(x, p) = (2, 2)$, so adding spatial
dimensions strengthens the linear contraction rate proportionally.

Two natural directions for further study are:
(i) a Steffensen-type acceleration in the symmetric case to upgrade
linear convergence to quadratic;
(ii) the asymmetric Jacobian and its rank-one decomposition when the
ratios $x_i$ differ.  Both are within reach by adapting the techniques
of~\cite{0pandrosion_pth} from 2D to $dD$.

\begin{thebibliography}{99}

\bibitem{0pandrosion_pth}
A formally verified Pandrosion--Steffensen scheme, $\texttt{0pandrosion\_pth.tex}$,
2025.

\end{thebibliography}

\end{document}
